% !TEX program = xelatex
% !TEX encoding = UTF-8
%
% 编译：将 WanPhys-dev/output/ 整夹复制到与本 .tex 同级目录下，再：
%   xelatex LBM流体仿真模块_报告.tex
%   xelatex LBM流体仿真模块_报告.tex
%
% 插图路径均为 output/...（相对 .tex 所在目录）。

\documentclass[UTF8,a4paper,11pt]{ctexart}

% ---------- 版式与字体 ----------
\usepackage{geometry}
\geometry{margin=2.5cm, headheight=14pt}

% ---------- 数学 ----------
\usepackage{amsmath,amssymb,bm,mathtools}

% ---------- 表格与列表 ----------
\usepackage{booktabs}
\usepackage{array}
\usepackage{longtable}
\usepackage{enumitem}
\setlist{nosep}

% ---------- 图表 ----------
\usepackage{graphicx}
\usepackage{float}
\usepackage{caption}
\usepackage{subcaption}
\captionsetup{font=small, labelfont=bf}

% ---------- 代码与终端输出 ----------
\usepackage{xcolor}
\usepackage{listings}
\lstset{
  basicstyle=\ttfamily\small,
  breaklines=true,
  frame=single,
  backgroundcolor=\color{gray!4},
  keywordstyle=\color{blue!70!black},
  commentstyle=\color{green!40!black},
  numbers=left,
  numberstyle=\tiny\color{gray},
  tabsize=2,
  showstringspaces=false,
}

% ---------- 超链接（放在导言区末尾） ----------
\usepackage{hyperref}
\hypersetup{
  colorlinks=true,
  linkcolor=blue,
  citecolor=blue,
  urlcolor=blue,
  pdftitle={基于 WanPhys 的 D3Q19 LBM 流体仿真模块},
  pdfauthor={课程项目组},
}

% ---------- 常用数学符号 ----------
\newcommand{\vecx}{\mathbf{x}}
\newcommand{\vecu}{\mathbf{u}}
\newcommand{\vece}{\mathbf{e}}
\newcommand{\feq}{f^{\mathrm{eq}}}
\newcommand{\dt}{\Delta t}
\newcommand{\dx}{\Delta x}
\newcommand{\cs}{c_s}
\newcommand{\Rey}{\mathrm{Re}}

\title{\textbf{基于 WanPhys 的\\D3Q19 格子玻尔兹曼流体仿真模块}\\[0.6em]
  \large 项目技术报告}
\author{WanPhys LBM 项目组}
\date{2026 年 7 月}

\begin{document}
\maketitle

\begin{abstract}
  本项目在 WanPhys \texttt{fluid\_grid} 框架内实现 D3Q19 单相流 LBM 模块，
  遵循与 \texttt{basic\_vortex}、\texttt{liquid}、\texttt{apic} 相同的
  Model / State / Solver / Domain 分层，继承 \texttt{FluidGridModelBase}、
  \texttt{FluidGridStateBase}、\texttt{FluidGridSolverBase} 及 Newton \texttt{Domain} 契约，
  在 GPU（Warp）上完成 BGK / TRT / MRT 碰撞、Pull 迁移、六面边界与格点固体障碍。
  示例通过 Newton 运行壳驱动 \texttt{FluidGridLbmDomain.step()}，后处理复用
  \texttt{fluid\_viewer} 与 VTK 管线。
  报告说明如何在 WanPhys 框架内扩展新算法族、核心实现、验证结果及已知局限。
\end{abstract}

\tableofcontents
\newpage

%======================================================================
\section{课题背景与目标}
%======================================================================

\subsection{背景}

流体模拟在图形学、机器人与游戏中应用广泛。传统欧拉网格法直接离散速度、压力等宏观量，
通过压力投影实现近似不可压缩；格子玻尔兹曼方法（LBM）则在规则格点上存储沿离散速度方向
传播的分布函数 $f_i$，通过\textbf{局部碰撞}与\textbf{邻格迁移}恢复宏观密度 $\rho$、
速度 $\vecu$ 与压力。

\subsection{作业要求摘要}

依据课程说明（\texttt{newton.txt}），基础目标包括：
\begin{itemize}
  \item 分布函数相关计算、BGK 碰撞、streaming 传播；
  \item 基本边界条件、简单外力项；
  \item 稳定性验证、示例场景与可视化。
\end{itemize}
可选扩展含 MRT/TRT 碰撞格式对比、复杂边界与固体耦合等。本组在基础 Pipeline 之上
实现了绕障流与 BGK/MRT 对比。

\subsection{预期成果与完成情况}

\begin{table}[H]
  \centering
  \caption{基础要求与实现对照}
  \begin{tabular}{@{}lll@{}}
    \toprule
    要求项 & 实现位置 & 状态 \\
    \midrule
    D3Q19 数据结构 & \texttt{lattice.py}, \texttt{state.py} & 已完成 \\
    BGK 碰撞 & \texttt{kernels.collide\_bgk} & 已完成 \\
    Pull streaming & \texttt{kernels.stream\_pull} & 已完成 \\
    边界条件 & \texttt{kernels.apply\_boundaries} & 已完成（简化格式） \\
    简单外力 & \texttt{kernels.update\_macro}（Guo 简化） & 接口已有，演示默认关闭 \\
    示例与可视化 & \texttt{examples/}, \texttt{cavity\_plot.py} 等 & 已完成 \\
    MRT / TRT（扩展） & \texttt{collide\_mrt}, \texttt{collide\_trt} & 已完成 \\
    绕障流 & \texttt{fluid\_grid\_lbm\_obstacle.py} 等 & 已完成 \\
    \bottomrule
  \end{tabular}
\end{table}

%======================================================================
\section{总体架构}
%======================================================================

\subsection{嵌入 WanPhys \texttt{fluid\_grid} 框架}

课程要求基于 Newton \texttt{release-1.1} / WanPhys 实现 LBM；\textbf{物理求解与数据结构深度落在 WanPhys 内}，
而非在引擎外另起一套脚本。本模块作为 \texttt{fluid\_grid/} 下的\textbf{第四条算法支线}
（与 \texttt{basic\_vortex}、\texttt{liquid}、\texttt{apic} 并列），路径为
\texttt{wanphys/\_src/fluid/fluid\_grid/lbm/}，并已从 \texttt{fluid\_grid/\_\_init\_\_.py} 导出
\texttt{FluidGridLbmModel} 等公开类型。

\begin{table}[H]
  \centering
  \caption{与 WanPhys \texttt{fluid\_grid} 抽象层的对接}
  \begin{tabular}{@{}llp{7.2cm}@{}}
    \toprule
    WanPhys 基类 / 契约 & LBM 实现 & 说明 \\
    \midrule
    \texttt{FluidGridModelBase} & \texttt{FluidGridLbmModel} & $\nu$, BC 类型, \texttt{collide\_impl} \\
    \texttt{FluidGridStateBase} & \texttt{FluidGridLbmState} & $f$, $F$, $\rho$, $\vecu$；复用 \texttt{solid\_phi} 烘焙 \\
    \texttt{FluidGridSolverBase} & \texttt{FluidGridLbmSolver} & 重写 \texttt{step()}，编排 LBM 五步 \\
    \texttt{Domain} & \texttt{FluidGridLbmDomain} & 双缓冲 \texttt{create\_state}/\texttt{step}，与 \texttt{liquid/domain} 同模式 \\
    \bottomrule
  \end{tabular}
\end{table}

LBM 以分布函数 $f_i$ 为主状态、通过碰撞松弛隐含不可压，\textbf{不适用}
\texttt{FluidGridMacSolverBase} 的 MAC 压力投影 Hook（对流$\to$散度$\to$泊松$\to$投影）。
这与 \texttt{liquid}/\texttt{apic} 虽继承 MAC 基类但\textbf{重写整段 \texttt{step()}}、
走粒子--网格管线的做法同类：\textbf{在 WanPhys 统一 Solver/Domain 接口下，实现算法专用的推进逻辑}，
而不是脱离框架的“外挂求解器”。详见 \texttt{fluid\_grid/ARCHITECTURE.md} 与 \texttt{DESIGN.md}。

\subsection{Newton 示例壳的角色}

Newton \textbf{不提供} LBM 内核；其职责是\textbf{运行已注册的 WanPhys Domain}：
\texttt{newton.examples.run()} 负责 CLI、帧循环与 \texttt{test\_final()} 钩子。
\texttt{fluid\_grid\_lbm\_*.py} 中实例化 \texttt{FluidGridLbmDomain}、配置 BC/障碍后调用
\texttt{domain.step()}——与 \texttt{basic\_vortex}、\texttt{liquid} 示例调用各自 Domain 的方式一致。
后处理（\texttt{cavity\_plot}、\texttt{vtk\_export}、\texttt{lbm\_flow}）读取 \texttt{state.v.numpy()}，
同样挂在 WanPhys 流体可视化子模块下。

\begin{table}[H]
  \centering
  \caption{职责划分}
  \begin{tabular}{@{}p{3.2cm}p{4.5cm}p{6cm}@{}}
    \toprule
    层次 & 归属 & 职责 \\
    \midrule
    算法 + 数据 & \textbf{WanPhys} \texttt{fluid\_grid/lbm} & Model/State/Solver/Domain、Warp kernels \\
    场景编排 & WanPhys \texttt{examples/} & 腔体/绕障 IC/BC，组装 Domain \\
    帧循环壳 & Newton \texttt{examples} & \texttt{--num-frames}、Viewer、验收回调 \\
    后处理 & WanPhys 绘图 / VTK & PNG、ParaView、可选 GL \\
    \bottomrule
  \end{tabular}
\end{table}

\subsection{核心类与单步流程}

模块包含 \texttt{FluidGridLbmModel}（静态参数）、
\texttt{FluidGridLbmState}（$f$, $F$, $\rho$, $\vecu$, \texttt{solid}）、
\texttt{FluidGridLbmSolver}（单步编排）、
\texttt{FluidGridLbmDomain}（双缓冲）。

单步顺序（已冻结）：
\begin{equation}
  \text{collide} \;\to\; \text{stream\_pull} \;\to\;
  \text{apply\_boundaries} \;\to\; \text{update\_macro} \;\to\; \text{swap}(f, F)
\end{equation}

选用 \textbf{Pull} 迁移：格点从上游 $x-\vece_i$ 读取分布，配合双缓冲避免 GPU 竞态，
并与参考实现 \texttt{taichi\_LBM3D} 的 \texttt{streaming1} 一致，便于对照。

%======================================================================
\section{核心算法}
%======================================================================

\subsection{D3Q19 与格子单位}

采用 D3Q19 模型（$Q=19$ 离散速度），平衡态（格子单位，$\cs^2=1/3$）：
\begin{equation}
  \feq_i(\rho, \vecu) = w_i \rho \left[
    1 + \frac{3\vece_i\cdot\vecu}{\cs^2}
    + \frac{9(\vece_i\cdot\vecu)^2}{2\cs^4}
    - \frac{3\|\vecu\|^2}{2\cs^2}
  \right].
\end{equation}

全程使用格子单位：$\dx=\dt=1$，运动黏度 $\nu$ 由模型传入，
松弛时间 $\tau = 3\nu + 0.5$，$\omega = 1/\tau$。
\textbf{不对照文献前禁止}用 \texttt{cell\_size} 换算物理速度。

\subsection{BGK 碰撞}

单松弛时间 BGK：
\begin{equation}
  f_i' = f_i - \omega\,(f_i - \feq_i(\rho, \vecu)).
\end{equation}
固体格点（\texttt{solid}$=1$）跳过碰撞，由迁移与 bounce-back 处理。

\subsection{TRT 与 MRT（扩展）}

\textbf{TRT}：对分布的正负对称部分使用不同松弛率 $\omega_\pm$，改善高 Re 稳定性。

\textbf{MRT}：将 $f$ 变换到矩空间，对各矩独立松弛后再变回速度空间。
本实现采用\textbf{简化松弛矩阵}（\texttt{mrt.py}）：
\begin{itemize}
  \item 守恒矩（$\rho$, $j_x$, $j_y$, $j_z$）：松弛率 $0$；
  \item 剪切相关 hydrodynamic 矩（4--15）：松弛率 $\omega$（与 BGK 相同）；
  \item ghost 矩（16--18）：松弛率 \texttt{mrt\_ghost\_s}$=1.0$。
\end{itemize}
因此在稳定工况下 MRT 宏观场与 BGK 接近；\textbf{主要优势}体现在高 Re / 粗网格下
BGK 发散时 MRT 仍可计算。MRT 碰撞核复杂度约 $O(19^3)$/格点，显著慢于 BGK。

\subsection{外力（Guo 简化）}

体积力配置于 \texttt{FluidGridLbmModel.force}，开关 \texttt{use\_guo\_force}。
仅在 \texttt{update\_macro} 中对宏观速度做修正：
\begin{equation}
  \vecu \leftarrow \frac{1}{\rho}\sum_i \vece_i f_i + \frac{\mathbf{F}}{2\rho},
\end{equation}
其中 $\mathbf{F}$ 为格子单位体积力；\textbf{碰撞步未写入完整 Guo 源项}。

当前演示算例（腔体、绕障、通道、BGK/MRT 对比）均设
\texttt{use\_guo\_force=False}，流动由\textbf{边界驱动}（顶盖速度、入口速度），
\textbf{未启用}体积外力。

%======================================================================
\section{边界条件}
%======================================================================

\subsection{六面分类框架}

计算域六个外表面各自独立配置类型码 \texttt{bc\_*}：
\begin{table}[H]
  \centering
  \caption{边界类型码}
  \begin{tabular}{@{}cll@{}}
    \toprule
    码 & 类型 & 实现 \\
    \midrule
    0 & 周期 & \texttt{stream\_pull} 内 wrap，BC 跳过 \\
    1 & 压力型 & $F_i = \feq_i(\rho_{\mathrm{bc}}, \vecu_{\mathrm{inner}})$ \\
    2 & 速度型 & $F_i = \feq_i(\rho=1, \vecu_{\mathrm{bc}})$ \\
    \bottomrule
  \end{tabular}
\end{table}

\textbf{说明}：当前为 taichi 参考的\textbf{平衡态填边界}，非完整 Zou--He 反推未知分布。

\subsection{顶盖驱动方腔}

坐标约定：\texttt{velocity[i,j,k]} 中 $i=x$, $j=y$, $k=z$；顶盖位于 $j=ny-1$
（字段 \texttt{bc\_y\_right}）。

\begin{lstlisting}[language=Python,caption={腔体边界配置}]
domain.solver.configure_cavity_walls()   # 六面 bc=2, u=0
domain.solver.set_lid_velocity(wp.vec3(U_lid, 0, 0))  # 顶盖
\end{lstlisting}

六面均为\textbf{速度型}，因物理边界给定的是壁面/顶盖速度（无滑移与匀速顶盖），
而非定压进出口。$z$ 向两端亦为 $u=0$，故为\textbf{3D 封闭盒}，非理想 2D 无限展向腔。

\subsection{通道绕障}

六面\textbf{异质}配置：$x=0$ 速度入口；$x=nx-1$ 压力出口；$y$ 壁面无滑移；$z$ 周期。
入口为\textbf{均匀来流}，无发展段与剖面入口。

\subsection{已知 BC 缺陷}

\begin{enumerate}
  \item \textbf{非完整 Zou--He}：19 个方向全部设为 $\feq$，角点 $\rho$ 误差偏大。
  \item \textbf{角点未专门处理}：棱边格点被多个面依次重写 $\feq$，
        顶盖--侧壁交点（如 $(0, ny-1, k)$）可能被顶盖速度覆盖侧壁 $u=0$，
        导致底角/顶角副涡偏弱或与 2D 文献偏差。
  \item \textbf{固体表面}：障碍用 \texttt{solid} + bounce-back，圆柱在格点上呈阶梯形。
\end{enumerate}

%======================================================================
\section{演示算例}
%======================================================================

\subsection{顶盖方腔}

脚本：\texttt{wanphys/examples/fluid\_grid\_lbm\_cavity.py}。
五壁固定 + 顶盖 $u_x=U_{\mathrm{lid}}$ 驱动主涡；输出中截面（$k=n_z/2$）流线 PNG 与 VTK。

\begin{figure}[H]
  \centering
  \begin{subfigure}[t]{0.48\textwidth}
    \centering
    \includegraphics[width=\textwidth]{output/cavity_speed.png}
    \caption{本仿真结果}
  \end{subfigure}\hfill
  \begin{subfigure}[t]{0.48\textwidth}
    \centering
    \includegraphics[width=\textwidth]{output/cavity_ref.png}
    \caption{参考图}
  \end{subfigure}
  \caption{顶盖驱动方腔速度流线（中截面 $k=n_z/2$）。
    参数：默认 $50^3$ 网格，$\nu=0.16667$，$U_{\mathrm{lid}}=0.1$（格子单位），BGK，300 帧（1500 步），$\Rey\approx 30$。}
  \label{fig:cavity-basic}
\end{figure}

\subsection{腔体绕障}

脚本：\texttt{fluid\_grid\_lbm\_obstacle.py}。
中央方柱，柱高 \texttt{mid / high / tall}；每种高度 6 张流线 PNG + 带 \texttt{solid} 的 VTK。
默认 $64^3$，$\nu=0.16667$，$U_{\mathrm{lid}}=0.1$，顶盖驱动。

\begin{figure}[H]
  \centering
  \begin{subfigure}[t]{0.32\textwidth}
    \centering
    \includegraphics[width=\textwidth]{output/obstacle/mid/stream_xy_z_mid.png}
    \caption{MID（柱居中）}
  \end{subfigure}\hfill
  \begin{subfigure}[t]{0.32\textwidth}
    \centering
    \includegraphics[width=\textwidth]{output/obstacle/high/stream_xy_z_mid.png}
    \caption{HIGH（底对齐，向上贴顶盖）}
  \end{subfigure}\hfill
  \begin{subfigure}[t]{0.32\textwidth}
    \centering
    \includegraphics[width=\textwidth]{output/obstacle/tall/stream_xy_z_mid.png}
    \caption{TALL（自底至顶）}
  \end{subfigure}
  \caption{腔体绕障：侧面 $x$--$y$ 速度流线（障碍 $z$ 向中间切片 \texttt{stream\_xy\_z\_mid}）。
    参数：$64^3$，$\nu=0.16667$，$U_{\mathrm{lid}}=0.1$，BGK，300 帧（1500 步）。}
  \label{fig:obstacle-heights}
\end{figure}

\subsection{通道绕障}

脚本：\texttt{fluid\_grid\_lbm\_channel\_obstacle.py}。
无穷高圆柱/方柱；俯视图尾迹与 $u_x,u_z$ 剖面。
默认 $240\times 80\times 80$，$\nu=0.012$，$u_{\mathrm{in}}=0.06$，障碍直径 $D=14$ 格，$\Rey\approx u_{\mathrm{in}} D/\nu \approx 70$。

\begin{figure}[H]
  \centering
  \begin{subfigure}[t]{0.48\textwidth}
    \centering
    \includegraphics[width=\textwidth]{output/channel_obstacle/cylinder/fig1_top_view_streamlines.png}
    \caption{本仿真（圆柱俯视）}
  \end{subfigure}\hfill
  \begin{subfigure}[t]{0.48\textwidth}
    \centering
    \includegraphics[width=\textwidth]{output/channel_obstacle/channel_ref.png}
    \caption{参考图}
  \end{subfigure}
  \caption{通道绕圆柱：俯视 $|u|$ 填色与 $(u_x,u_z)$ 流线。
    格子单位：$240\times 80\times 80$，$\nu=0.012$，$u_{\mathrm{in}}=0.06$，$D=14$，300 帧。}
  \label{fig:channel-ref}
\end{figure}

\begin{figure}[H]
  \centering
  \includegraphics[width=0.72\textwidth]{output/channel_obstacle/cylinder/fig2_slice_ux_uz.png}
  \caption{通道绕圆柱：$y=n_y/2$ 剖面 $u_x$、$u_z$ 曲线（参数同图~\ref{fig:channel-ref}）。}
  \label{fig:channel-cylinder}
\end{figure}

\begin{figure}[H]
  \centering
  \begin{subfigure}[t]{0.48\textwidth}
    \centering
    \includegraphics[width=\textwidth]{output/channel_obstacle/box/fig1_top_view_streamlines.png}
    \caption{方柱：俯视流线}
  \end{subfigure}\hfill
  \begin{subfigure}[t]{0.48\textwidth}
    \centering
    \includegraphics[width=\textwidth]{output/channel_obstacle/box/fig2_slice_ux_uz.png}
    \caption{方柱：$u_x$、$u_z$ 剖面曲线}
  \end{subfigure}
  \caption{通道绕无穷高方柱（网格与来流参数同图~\ref{fig:channel-cylinder}）。}
  \label{fig:channel-box}
\end{figure}

\subsection{BGK / MRT 对比}

脚本：\texttt{fluid\_grid\_lbm\_cavity.py --compare}；参数预设见 \texttt{cavity\_presets.py}。
雷诺数估计：$\Rey \approx U_{\mathrm{lid}}\,N/\nu$（$N$ 为网格边长格点数）。

\begin{table}[H]
  \centering
  \caption{对比预设}
  \begin{tabular}{@{}lrrrrr@{}}
    \toprule
    profile & 网格 & $\nu$ & $U_{\mathrm{lid}}$ & 步数 & $\Rey$ \\
    \midrule
    subtle  & $64^3$  & 0.16667 & 0.10 &  800 & $\sim 38$ \\
    stress  & $64^3$  & 0.025   & 0.16 & 1500 & $\sim 410$ \\
    contrast& $128^3$ & 0.01024 & 0.42 & 2500 & $\sim 5250$ \\
    coarse  & $32^3$  & 0.00256 & 0.20 & 2500 & $\sim 2500$ \\
    \bottomrule
  \end{tabular}
\end{table}

\textbf{coarse} 档为答辩推荐：BGK 约在 1600--1700 步后出现 NaN；MRT 2500 步仍稳定。
\texttt{--compare} 依次跑 BGK 与 MRT，并生成 \texttt{cavity\_diff\_*.png} 差分图
（BGK 发散时可能无 BGK 的 NPZ，属预期）。

\begin{figure}[H]
  \centering
  \begin{subfigure}[t]{0.48\textwidth}
    \centering
    \includegraphics[width=\textwidth]{output/cavity_compare/subtle/cavity_bgk_speed.png}
    \caption{subtle / BGK}
  \end{subfigure}\hfill
  \begin{subfigure}[t]{0.48\textwidth}
    \centering
    \includegraphics[width=\textwidth]{output/cavity_compare/subtle/cavity_mrt_speed.png}
    \caption{subtle / MRT}
  \end{subfigure}
  \caption{低 Re 对比（subtle）：$64^3$，$\nu=0.16667$，$U_{\mathrm{lid}}=0.1$，800 步，$\Rey\approx 38$；
    BGK 与 MRT 宏观场应几乎一致。}
  \label{fig:compare-subtle}
\end{figure}

\begin{figure}[H]
  \centering
  \begin{subfigure}[t]{0.48\textwidth}
    \centering
    \includegraphics[width=\textwidth]{output/cavity_compare/stress/cavity_bgk_speed.png}
    \caption{stress / BGK}
  \end{subfigure}\hfill
  \begin{subfigure}[t]{0.48\textwidth}
    \centering
    \includegraphics[width=\textwidth]{output/cavity_compare/stress/cavity_mrt_speed.png}
    \caption{stress / MRT}
  \end{subfigure}
  \caption{中 Re 对比（stress）：$64^3$，$\nu=0.025$，$U_{\mathrm{lid}}=0.16$，1500 步，$\Rey\approx 410$。}
  \label{fig:compare-stress}
\end{figure}

\begin{figure}[H]
  \centering
  \begin{subfigure}[t]{0.48\textwidth}
    \centering
    \includegraphics[width=\textwidth]{output/cavity_compare/contrast/cavity_bgk_speed.png}
    \caption{contrast / BGK}
  \end{subfigure}\hfill
  \begin{subfigure}[t]{0.48\textwidth}
    \centering
    \includegraphics[width=\textwidth]{output/cavity_compare/contrast/cavity_mrt_speed.png}
    \caption{contrast / MRT}
  \end{subfigure}
  \caption{高 Re 细网格（contrast）：$128^3$，$\nu=0.01024$，$U_{\mathrm{lid}}=0.42$，2500 步，$\Rey\approx 5250$。}
  \label{fig:compare-contrast}
\end{figure}

\begin{figure}[H]
  \centering
  \includegraphics[width=0.72\textwidth]{output/cavity_compare/coarse/cavity_mrt_speed.png}
  \caption{粗网格高 Re（coarse）：$32^3$，$\nu=0.00256$，$U_{\mathrm{lid}}=0.2$，2500 步，$\Rey\approx 2500$。
    同参数下 BGK 约 1600+ 步发散（本目录无 BGK 速度图）；MRT 仍稳定，见 \texttt{cavity\_mrt\_speed.png}。}
  \label{fig:compare-coarse}
\end{figure}

%======================================================================
\section{验证与测试}
%======================================================================

\begin{table}[H]
  \centering
  \caption{主要测试}
  \begin{tabular}{@{}ll@{}}
    \toprule
    测试 & 通过标准 \\
    \midrule
    \texttt{test\_lbm\_rest.py} & 16$^3\times$100 步，$\max\|\vecu\|<10^{-10}$，$\max|\rho-1|<10^{-6}$ \\
    \texttt{test\_cavity\_smoke.py} & 32$^3\times$500 步，流动建立，无 NaN \\
    \texttt{test\_obstacle\_smoke.py} & 绕障不穿透、柱后有 wake \\
    \bottomrule
  \end{tabular}
\end{table}

终端每 20 帧打印 \texttt{max|u|}、\texttt{near\_lid\_ux}、\texttt{max|rho-1|} 供监控。
\texttt{test\_final()} 以 NaN/Inf 判定发散；有限但异常的大速度仍可能显示 \texttt{Cavity OK}。

%======================================================================
\section{不足与后续工作}
%======================================================================

\begin{itemize}
  \item 边界：完整 Zou--He、角点格式、通道入口发展段。
  \item 几何：曲边界 / IBM，改善圆柱阶梯离散。
  \item MRT：文献完整 $S$ 矩阵；性能优化。
  \item 验证：与 Ghia 等基准的涡心位置定量对比。
  \item 外力：碰撞步补全 Guo 源项。
\end{itemize}

%======================================================================
\section{结论}
%======================================================================

本项目在 WanPhys \texttt{fluid\_grid} 框架内实现了完整的 D3Q19 LBM 算法族：
遵循 Model/State/Solver/Domain 分层与 \texttt{FluidGridSolverBase} 契约，GPU 并行、Pull 迁移、
六面 BC 与格点障碍均落在 WanPhys 模块内；Newton 仅作示例运行壳。
基础 BGK 满足课程要求；TRT/MRT 与绕障流体现框架内可扩展性。
BGK/MRT 对比表明：低 Re 下两者一致；高 Re 粗网格下 MRT 稳定性优于 BGK。
后续可在 WanPhys 框架内继续完善边界、曲边界耦合与定量验证。

%======================================================================
\appendix
\section{常用编译与运行命令}
%======================================================================

\begin{lstlisting}[language=bash,caption={环境}]
cd WanPhys-dev
.venv_lbm/Scripts/Activate.ps1   # Windows
\end{lstlisting}

\begin{lstlisting}[language=bash,caption={顶盖方腔}]
python wanphys/examples/fluid_grid_lbm_cavity.py \
  --viewer null --num-frames 300 --test \
  --output-dir output/video/cavity
\end{lstlisting}

\begin{lstlisting}[language=bash,caption={BGK/MRT 对比（答辩推荐 coarse）}]
python wanphys/examples/fluid_grid_lbm_cavity.py \
  --viewer null --test --profile coarse --compare \
  --output-dir output/cavity_compare/coarse
\end{lstlisting}

\begin{lstlisting}[language=bash,caption={绕障}]
python wanphys/examples/fluid_grid_lbm_obstacle.py \
  --viewer null --num-frames 300 --test --obstacle-height all
python wanphys/examples/fluid_grid_lbm_channel_obstacle.py \
  --viewer null --num-frames 300 --test --obstacle-mode both
\end{lstlisting}

\section{参考文献}
\begin{enumerate}[label={[\arabic*]}]
  \item Newton Physics Engine: \url{https://github.com/newton-physics/newton}
  \item Kr\"uger et al., \textit{The Lattice Boltzmann Method: Principles and Practice}.
  \item taichi\_LBM3D: \url{https://github.com/yjhp1016/taichi_LBM3D}
  \item 本项目 \texttt{DESIGN.md}, \texttt{LBM演示算例\_详细说明.md}
\end{enumerate}

\end{document}
